\documentclass[11pt,a4paper]{article}

%% PACHETE MATEMATICE
\usepackage{amsmath,amssymb,amsfonts}
\usepackage{amsthm}
\usepackage{mathpazo}
\usepackage{eulervm}
\usepackage{enumerate}
\usepackage{color}
\usepackage{graphicx}
\usepackage[all]{xy}
\usepackage{xcolor}
\usepackage{framed}
\usepackage[unicode]{hyperref}
\setcounter{MaxMatrixCols}{20}
\usepackage{multicol}
%\usepackage[romanian]{babel}


%% ASPECT
\linespread{1.05}
\usepackage[T1]{fontenc}
\usepackage[utf8x]{inputenc}
\usepackage[margin=2cm]{geometry}
% \usepackage[color]{showkeys}
\usepackage{anyfontsize}
\usepackage{titlesec}
\titleformat*{\section}{\large\bfseries}

\usepackage{fancyhdr}
\pagestyle{fancy}
\rhead{\small \color{gray} Notițe de Adrian Manea}
\lhead{\small \color{gray} ACS-M1, 2017---2018, M. Olteanu}


\usepackage[autostyle = false, style = german]{csquotes}
\usepackage[romanian]{babel}
\newcommand\qq{\enquote}

%% COMENZILE MELE
\renewcommand*{\proofname}{Dem.:}
\newcommand\bb{\mathbb}
\newcommand\dr{\mathrm}
\newcommand\kal{\mathcal}
\newcommand\fk{\mathfrak}
\newcommand\op{\oplus}
\newcommand\ot{\otimes}
\newcommand\ds{\displaystyle}
\newcommand\id{\indent}
\newcommand\nid{\noindent}
\newcommand\seq{\subseteq}
\newcommand\sneq{\subsetneq}
\newcommand\speq{\supseteq}
\newcommand\spneq{\supsetneq}
\newcommand\rar{\rightarrow}
\newcommand\Rar{\Rightarrow}
\newcommand\xrar{\xrightarrow}
\newcommand\lar{\leftarrow}
\newcommand\Lar{\Leftarrow}
\newcommand\xlar{\xleftarrow}
\newcommand\lrar{\leftrightarrow}
\newcommand\Lrar{\Leftrightarrow}
\newcommand\ideal{\trianglelefteq}
\newcommand\ai{\text{ a.\^{i}. }}
\newcommand\sk{(\textit{Schi\c{t}\u{a}}) }
\newcommand\rhup{\rightharpoonup}
\newcommand\rhdn{\rightharpoondown}
\newcommand\lhup{\leftharpoonup}
\newcommand\lhdn{\leftharpoondown}
\newcommand\vid{\emptyset}
\newcommand\wh{\widehat}
\newcommand\ol{\overline}
\newcommand\NN{\mathbb{N}}
\newcommand\RR{\mathbb{R}}
\newcommand\ZZ{\mathbb{Z}}
\newcommand\QQ{\mathbb{Q}}
\newcommand\CC{\mathbb{C}}

%% STILURILE MELE
\newtheoremstyle{dotless}{}{}{\itshape}{}{\bfseries}{:}{ }{}

\theoremstyle{dotless}
\newtheorem{theorem}{Teorem\u{a}}[section]
\newtheorem{proposition}{Propozi\c{t}ie}[section]
\newtheorem{lemma}{Lem\u{a}}[section]
\newtheorem{corollary}{Corolar}[section]
\newtheorem{exercise}{Exerci\c{t}iu}

\newtheoremstyle{dotlessdr}{}{}{}{}{\bfseries}{:}{ }{}
\theoremstyle{dotlessdr}
\newtheorem{example}{Exemplu}[section]
\newtheorem{definition}{Defini\c{t}ie}[section]
\newtheorem{remark}{Observa\c{t}ie}[section]




\begin{document}
% \definecolor{labelkey}{RGB}{222,0,0}

\begin{center}
  {\large\textbf{Seminar 1 \\ Șiruri și serii de numere reale}}
\end{center}

\section{Serii --- Noțiuni generale}

\begin{definition}\label{def:sir}
  Se numește \emph{șir de numere reale (complexe)} orice funcție de forma
  $f : \mathbb{N} \to \mathbb{R}$ (respectiv $\mathbb{C}$).

  În această notație, se obișnuiește ca în loc de imaginea numărului
  natural $n$ să se scrie sub formă de indice, i.e.\ $x_n = f(n)$, iar
  șirul în ansamblu să se noteze cu $(x_n)_{n \in \NN}$.
\end{definition}

\begin{definition}\label{def:serie}
  Fie $(x_n)_n$ un șir de numere reale și fie seria $\sum_n x_n$, definită
  de șirul sumelor parțiale $s_n = \ds\sum_{k=1}^n u_k$.

  Seria se numește \emph{convergentă} dacă șirul $s_n$ este convergent, iar
  limita șirului se numește \emph{suma seriei}. În caz contrar, seria se numește
  \emph{divergentă}, adică șirul sumelor parțiale nu are limită sau aceasta este
  infinită.
\end{definition}

Cîteva proprietăți, care ne ajută să lucrăm cu operații cu serii:
\begin{enumerate}[(1)]
  \item Dacă într-o serie schimbăm ordinea unui număr finit de termeni, obținem
    o serie nouă, care are aceeași natură. Dacă există, suma seriei nu se schimbă.
  \item Dacă eliminăm un număr finit de termeni dintr-o serie, se obține o nouă serie
    cu aceeași natură. Suma, dacă există, poate să se schimbe.
  \item Fie seria $\sum_n x_n$ convergentă. Atunci pentru orice șir $(k_n)$ crescător
    și divergent, cu numere naturale, seria obținută prin reordonare:
    \[
      (x_1 + x_2 + \dots + x_{k_1}) + (x_{k_1 + 1} + \dots + x_{k_2}) + \dots
    \]
    este de asemenea convergentă și are aceeași sumă. Dacă seria inițială are sumă
    infinită, atunci și seria de mai sus are sumă infinită.

    Dacă există un șir $(k_n)$ ca mai sus, care să facă astfel încît noua serie
    reordonată să devină divergentă, atunci și seria inițială este divergentă.
  \item Dacă o serie este convergentă, atunci ea are șirul sumelor parțiale mărginit.
  \item Fie $\sum_{n \geq 1} x_n$ o serie convergentă și $\sum_{n \geq p} x_n$ seria
    obținută prin eliminarea primilor $p$ termeni. Suma seriei translatate se notează
    $R_p$ și se numește \emph{restul de ordin $p$} al seriei inițiale.

    Se poate observa faptul că resturile unei serii convergente formează un șir
    convergent către zero.
  \item Seriile $\sum_n x_n$ și $\sum_n \alpha x_n$, cu $\alpha \in \RR^\ast$ au
    aceeași natură.
  \item Dacă seria $\sum_n x_n$ este convergentă, atunci șirul $(x_n)$ al termenilor
    săi este convergent către zero. Reciproca este falsă $\Big( \sum_n \frac{1}{n} \Big)$.

    Proprietatea aceasta oferă, însă, \textbf{condiția necesară de convergență}, anume:
  \item Dacă șirul termenilor unei serii nu este convergent către zero, seria este
    divergentă.
\end{enumerate}

%%%%%%%%%%%%%%%%%%%%%%%%%%%%%%%%%%%%%%%%%%%%%%%%%%%%%%%%%%%%%%%%%%%%%%

\section{Serii cu termeni pozitivi}

Cîteva proprietăți simple care rezultă din particularizarea celor de mai sus:

\begin{enumerate}[(1)]
  \item Șirul sumelor parțiale ale unei serii cu termeni pozitivi este strict
    crescător;
  \item O serie cu termeni pozitivi are întotdeauna sumă (finită sau nu);
  \item O serie cu termeni pozitivi este convergentă dacă și numai dacă șirul
    sumelor parțiale este mărginit;
  \item Dacă $\sum_n x_n$ este o serie cu termeni pozitivi, atunci ea are aceeași
    natură cu seria obținută prin reordonare după un șir de numere naturale (cf.\
    proprietății (3) de mai sus). 
\end{enumerate}

Pentru a stabili natura unei serii cu termeni pozitivi, se pot folosi următoarele
criterii:

\begin{enumerate}[(1)]
  \item \textbf{Primul criteriu de comparație:} O serie care are termeni mai mari
    (doi cîte doi) decît o serie divergentă este divergentă. O serie care are termeni
    mai mici (doi cîte doi) decît o serie convergentă este convergentă.
  \item \textbf{Al doilea criteriu de comparație:} Fie $\sum_n x_n$ și $\sum_n y_n$
    două serii cu termeni pozitivi. Presupunem că $\frac{x_{n+1}}{x_n} \leq \frac{y_{n+1}}{y_n}$.
    Atunci:
    \begin{itemize}
      \item Dacă seria $\sum_n y_n$ este convergentă, atunci și seria $\sum_n x_n$ este
        convergentă;
      \item Dacă seria $\sum_n x_n$ este divergentă, atunci și seria $\sum_n y_n$ este
        divergentă.
    \end{itemize}
  \item \textbf{Al treilea criteriu de comparație:} Fie $\sum_n x_n$ și $\sum_n y_n$ serii
    cu termeni pozitivi, astfel încît $\ds\lim_{n \to \infty} \frac{x_n}{y_n} = \ell$.
    \begin{itemize}
      \item Dacă $0 < \ell < \infty$, atunci cele două serii au aceeași natură;
      \item Dacă $\ell = 0$, iar seria $\sum_n y_n$ este convergentă, atunci și seria
        $\sum_n x_n$ este convergentă;
      \item Dacă $\ell = +\infty$, iar seria $\sum_n y_n$ este divergentă, atunci seria
        $\sum_n x_n$ este divergentă.
    \end{itemize}
  \item \textbf{Criteriul radical (Cauchy)}: Fie $\sum_n x_n$ o serie cu termeni pozitivi,
    astfel încît $\ds\lim_{n \to \infty} \sqrt[n]{x_n} = \ell$. Atunci:
    \begin{itemize}
      \item Dacă $\ell < 1$, seria este convergentă;
      \item Dacă $\ell > 1$, seria este divergentă;
      \item Dacă $\ell = 1$, criteriul nu decide.
    \end{itemize}
  \item \textbf{Criteriul raportului (D'Alembert)}: Fie $\sum_n x_n$ o serie cu termeni
    pozitivi, astfel încît $\ds\lim_{n \to \infty} \frac{x_{n+1}}{x_n} = \ell$. Atunci:
    \begin{itemize}
      \item Dacă $\ell < 1$, seria este convergentă;
      \item Dacă $\ell > 1$, seria este divergentă;
      \item Dacă $\ell = 1$, criteriul nu decide.
    \end{itemize}
  \item \textbf{Criteriul lui Raabe-Duhamel:} Fie $\sum_n x_n$ o serie cu termeni
    pozitivi și notăm:
    \[
      \ell = \lim_{n \to \infty} n \Big( \frac{x_n}{x_{n+1}} - 1 \Big).
    \]
    Atunci:
    \begin{itemize}
      \item Dacă $\ell < 1$, seria este divergentă;
      \item Dacă $\ell > 1$, seria este convergentă;
      \item Dacă $\ell = 1$, criteriul nu decide.
    \end{itemize}
  \item \textbf{Criteriul logaritmic:} Fie $\sum_n x_n$ o serie cu termeni pozitivi
    și presupunem că există limita:
    \[
      \ell \stackrel{\text{not.}}{=\joinrel=} \lim_{n \to \infty} %
    \frac{\ln \frac{1}{x_n}}{\ln n}.
    \]
    Atunci:
    \begin{itemize}
      \item Dacă $\ell < 1$, seria este divergentă;
      \item Dacă $\ell > 1$, seria este convergentă;
      \item Dacă $\ell = 1$, criteriul nu decide.
    \end{itemize}
  \item \textbf{Criteriul integral:} Fie $f: (0, \infty) \to [0, \infty)$ o funcție
    crescătoare și șirul $a_n = \int_1^n f(t) dt$. Atunci seria $\sum_n f(n)$ este
    convergentă dacă și numai dacă șirul $(a_n)_n$ este convergent.
  \item \textbf{Criteriul condensării:} Fie $x_n \geq x_{n + 1} \geq 0, \forall n \in \NN$.
    Atunci seriile $\sum_n x_n$ și $\sum_n 2^n x_{2^n}$ au aceeași natură.
\end{enumerate}

%%%%%%%%%%%%%%%%%%%%%%%%%%%%%%%%%%%%%%%%%%%%%%%%%%%%%%%%%%%%%%%%%%%%%%

\section{Convergență absolută și semiconvergență}

Atunci cînd nu lucrăm cu serii care au doar termeni pozitivi, convergența trebuie
studiată ceva mai atent. De fapt, ea se subdivide în cel puțin două tipuri, pe
care le vom prezenta.

Însă, cel mai simplu exemplu de serie care nu are toți termenii pozitivi este
\emph{seria alternată:}

\begin{definition}\label{def:serie-alt}
  O serie $\sum_n x_n$ se numește \emph{alternată} dacă produsul
  $x_n \cdot x_{n+1} < 0, \forall n \in \NN$.
\end{definition}

Pentru astfel de serii, avem \textbf{criteriul lui Leibniz} care ne arată
convergența: Fie $\sum_n (-1)^{n+1} u_n$ o serie alternată. Dacă șirul
$(u_n)_n$ este descrescător și converge către $0$, atunci seria este
convergentă.

Ajungem, în fine, la noțiunile anunțate.

\begin{definition}\label{def:abs-conv}
  O serie $\sum_n x_n$ se numește \emph{absolut convergentă} dacă seria
  modulelor $\sum_n |x_n|$ este convergentă.
\end{definition}

Următoarele proprietăți sînt evidente:
\begin{itemize}
  \item Orice serie absolut convergentă este convergentă;
  \item Prin orice permutare a termenilor unei serii absolut convergente se
    obține una tot absolut convergentă și cu aceeași sumă.
\end{itemize}

\begin{remark}\label{rk:abs-conv-poz}
  Deoarece știm că pentru o serie absolut convergentă, proprietatea de
  convergență se păstrează și cînd trecem la module, putem să aplicăm pentru
  seriile absolut convergente criteriile din cazul seriilor cu termeni pozitivi.
\end{remark}


\begin{definition}\label{def:semiconv}
  O serie se numește \emph{semiconvergentă} dacă are seria modulelor divergentă.
\end{definition}

Așadar, dacă seria este absolut convergentă, ne putem reduce la cazul seriilor
cu termeni pozitivi, prin aplicarea modulului. În caz contrar (cel de semiconvergență),
avem la dispoziție următoarele criterii:

\begin{enumerate}[(1)]
  \item \textbf{Criteriul general (Cauchy)}: O serie $\sum_n x_n$ este convergentă
    dacă și numai dacă pentru orice $\varepsilon > 0$, există $N_\varepsilon \in \NN$
    astfel încît, pentru orice $n \geq N_\varepsilon$, are loc:
    \[
      |x_{n+1} + x_{n+2} + \dots + x_{n + p} | < \varepsilon, \forall p \geq 1.
    \]
    
  \item \textbf{Criteriul lui Abel}: Dacă seria $\sum_n x_n$ se poate scrie sub forma
    $\sum_n \alpha_n y_n$, cu $(\alpha_n)$ un șir monoton și mărginit, iar seria
    $\sum_n y_n$ este convergentă, atunci și seria inițială este convergentă.

  \item \textbf{Criteriul lui Dirichlet}: Dacă seria $\sum_n x_n$ se poate scrie
    sub forma $\sum_n \alpha_n y_n$, în care $(\alpha_n)$ este un șir monoton care
    tinde către zero, iar șirul de termen general $V_n = v_1 + \dots + v_n$ este
    mărginit, atunci seria inițială este convergentă.
\end{enumerate}

Mai avem următoarea noțiune:
\begin{definition}\label{def:conv-cond}
  O serie convergentă $\sum_n x_n$ se numește \emph{necondiționat convergentă} dacă
  pentru orice permutare $\sigma : \NN \to \NN$, seria $\sum_n x_{\sigma(n)}$ este
  convergentă. În caz contrar, seria se numește \emph{condiționat convergentă}.
\end{definition}

Următorul rezultat este evident:
\begin{theorem}[Dirichlet]\label{thm:abs-necond}
  Orice serie absolut convergentă este necondiționat convergentă.
\end{theorem}

În schimb, teorema care urmează poate fi surprinzătoare:
\begin{theorem}[Riemann]\label{thm:riemann-reord}
  Fiind dată o serie convergentă, dar nu absolut convergentă și un număr real $S$,
  există o permutare a termenilor seriei astfel încît noua serie obținută să aibă
  suma exact $S$.
\end{theorem}

%%%%%%%%%%%%%%%%%%%%%%%%%%%%%%%%%%%%%%%%%%%%%%%%%%%%%%%%%%%%%%%%%%%%%%

\section{Operații cu serii}

În unele situații, dacă o serie este mai complicată, putem să o descompunem
în două sau mai multe serii. Apoi, problema convergenței se reduce la studiul
componentelor, iar rezultatul final se formulează ținînd seama de operațiile
permise între serii (convergente).

Operațiile elementare cu serii sînt:
\begin{itemize}
  \item \emph{suma și diferența}, care se realizează termen cu termen, adică:
    \[
      \sum_n x_n + \sum_n y_n = \sum_n (x_n + y_n);
    \]
  \item \emph{produsul}, care se mai numește și \emph{convoluție}, deoarece nu
    se realizează termen cu termen, ci avem:
    \[
      \sum_n x_n \cdot \sum_n y_n = \sum_n \Big( \sum_j x_j y_{n - 1 - j} \Big).
    \]
\end{itemize}

Cu aceste operații, rezultatele principale care ne arată utilitatea lor sînt:

\begin{enumerate}[(1)]
  \item Suma și diferența a două serii convergente este o serie convergentă. Suma
    seriei rezultate se calculează corespunzător, i.e.\ ca suma sau diferența
    sumelor seriilor implicate.
  \item \textbf{Abel:} Suma unei serii produs se calculează ca produsul sumelor
    seriilor-factori.
  \item \textbf{Mertens:} Dacă seriile $\sum_n u_n, \sum_n v_n$ sînt convergente și
    cel puțin una dintre ele este absolut convergentă, atunci seria produs
    $\sum_n w_n$ este convergentă și avem $U \cdot V = W$, unde am notat cu majuscule
    sumele seriilor corespunzătoare.
  \item \textbf{Cauchy:} Dacă seriile factori sînt absolut convergente, atunci și seria
    produs este absolut convergentă.
\end{enumerate}


%%%%%%%%%%%%%%%%%%%%%%%%%%%%%%%%%%%%%%%%%%%%%%%%%%%%%%%%%%%%%%%%%%%%%%

\section{Aproximarea sumelor seriilor convergente}

Putem calcula numeric sumele unor serii convergente, cu anumite aproximații,
chiar dacă suma finală nu poate fi calculată în întregime (sau poate rezultatul
exact nu este de interes). În acest sens, avem următoarele două rezultate:

\begin{theorem}[Aproximarea sumelor seriilor cu termeni pozitivi]\label{thm:aprox-poz}
  
  Fie $x_n \geq 0$ și $k \geq 0$, astfel încît $\frac{x_{n+1}}{x_n} < k < 1, \forall n \in \NN$.
  Dacă $S$ este suma seriei convergente $\sum_n x_n$, iar $s_n$ este suma primilor $n + 1$
  termeni, atunci avem următoarea aproximație:
  \[
    | S - s_n | < \frac{k}{1 - k} x_n.
  \]
\end{theorem}

\begin{theorem}[Aproximarea sumelor seriilor alternate]\label{thm:aprox-alt}
  Fie $\sum_n (-1)^n x_n$ o serie alternată convergentă și fie $S$ suma sa. Dacă $S_n$ este suma
  primilor $n + 1$ termeni, atunci avem:
  \[
    | S - S_n | \leq x_{n + 1}.
  \]
\end{theorem}

În esență, ambele teoreme arată că \emph{eroarea aproximației are ordinul de mărime %
al primului termen neglijat}.

%%%%%%%%%%%%%%%%%%%%%%%%%%%%%%%%%%%%%%%%%%%%%%%%%%%%%%%%%%%%%%%%%%%%%%

\section{Exemple și exerciții}

Să studiem două serii care au fost deja întîlnite:

\textbf{Seria geometrică:} Fie $a \in \RR$ și $q \in \RR$, pe care îl numim \emph{rație}.
Considerăm seria geometrică $\sum_n aq^n$. După cum știm, suma parțială de rang $n$
a seriei este:
\[
  s_n = a + aq + \dots + aq^{n-1} = %
  \begin{cases}
    a \cdot \frac{1 - q^n}{1 - q}  & q \neq 1 \\
    na, & q = 1.
  \end{cases}
\]

În ce privește convergența, rezultă că dacă $|q| < 1$, atunci avem:
\[
  \lim_{n \to \infty} s_n = \frac{a}{1 - q},
\]
deci seria este convergentă și are suma chiar $\frac{a}{1 - q}$.

În caz contrar, dacă $|q| \geq 1$, se poate verifica folosind criteriul necesar că seria
geometrică este divergentă. \\

Un caz mai general îl constituie \textbf{seria armonică generalizată}, definită prin
$\sum_n \frac{1}{n^\alpha}$, cu $\alpha \in \RR$.

Dacă $\alpha \leq 0$, atunci termenul general al seriei nu converge către zero, deci din
criteriul necesar, obținem că seria este divergentă.

Dacă $\alpha > 0$, atunci termenii seriei formează un șir descrescător de numere pozitive.
Se poate constata, în acest caz, că seria este convergentă pentru $\alpha > 1$ și divergentă
pentru $\alpha \leq 1$.

Pentru $\alpha = 1$, seria se numește simplu \emph{seria armonică}.

\vspace{1cm}

1. Studiați natura următoarelor serii cu termeni pozitivi, definite de șirul $x_n$:
\begin{multicols}{2}
  \begin{enumerate}[(a)]
    \item $x_n = \Big( \frac{3n}{3n + 1} \Big)^n$ (D, necesar);
    \item $x_n = \frac{1}{n!}$ (C, comparație);
    \item $x_n = \frac{1}{n \sqrt{n + 1}}$ (C, comparație);
    \item $x_n = \frac{n!}{n^{2n}}$ (C, D'Alembert);
    \item $x_n = n! \cdot \Big( \frac{a}{n} \Big)^n, a \in \RR$ (discuție, raport);
    \item $x_n = \frac{(n!)^2}{(2n)!}$ (C, raport);
    \item $x_n = \frac{n!}{(a + 1)(a + 2) \cdots (a + n)}, a > -1$ (discuție, Raabe);
    \item $x_n = \frac{1 \cdot 3 \cdot 5 \cdots (2n - 1)}{2 \cdot 4 \cdot 6 \cdots 2n}$ (D, Raabe);
    \item $x_n = \Big( 1 - \frac{3 \ln n}{2n} \Big)^n$ (C, logaritmic);
    \item $x_n = \Big( \frac{n + 1}{3n + 1} \Big)^n$ (C, rădăcină);
    \item $x_n = \Big( \frac{an + 1}{bn + 1} \Big)^n, a, b > 0$ (discuție, rădăcină);
    \item $x_n = \Big( \frac{1}{\ln n} \Big)^{\ln(\ln n)}$ (D, logaritmic);
    \item $x_n = \frac{1}{n \ln n}$ (D, integral);
    \item $x_n = \frac{1}{n \ln^2 n}$ (C, integral);
    \item $x_n = \frac{\sqrt[n]{(n + 1) (n+2) \cdots (n + n)}}{n}$ (D, necesar + rădăcină);
    \item $x_n = \frac{1}{7^n + 3^n}$ (C, comparație);
    \item $x_n = 3^n \sin \frac{\pi}{5^n}$ (C, comparație $\sin x < x$);
    \item $x_n = \frac{1}{n \sqrt[n]{n}}$ (D, comparație 3 $\sim \frac{1}{n}$);
    \item $x_n = a^{\ln n}, a > 0$ (discuție, Raabe).
  \end{enumerate}
\end{multicols}

\vspace{1cm}

2. Studiați natura următoarelor serii, definite de șirul $x_n$:
\begin{enumerate}[(a)]
  \item $x_n = \frac{(-1)^n}{n}$ (C, Leibniz, nu AC);
  \item $x_n = (-1)^n$ (D, necesar);
  \item $x_n = \frac{1}{n + i}$ (D, descompune, una divergentă);
  \item $x_n = \frac{\cos n \alpha}{n^2}, \alpha \in \RR$ (AC, comparație);
\end{enumerate}

\vspace{1cm}

3. Dați exemplu de serii divergente a căror sumă este o serie convergentă.
($(-1)^n$ și $(-1)^{n + 1}$).

\vspace{1cm}

4. Să se aproximeze cu o eroare mai mică decît $\varepsilon$ sumele
seriilor definite de șirul $x_n$:
\begin{enumerate}[(a)]
  \item $x_n = \frac{(-1)^n}{n!}, \varepsilon = 10^{-3}$. ($x_n < \varepsilon \Rightarrow n = 7$,
    deci calculăm $S_7$);
  \item $x_n = \frac{(-1)^n}{n^3 \sqrt{n}}, \varepsilon = 10^{-2}$ ($n = 4$);
  \item $x_n = \frac{1}{n! 2^n}$ ($\frac{x_{n + 1}}{x_n} < \frac{1}{6} = k, \forall n \geq 2$,
    deci $S - s_n \leq x_n \frac{k}{1 - k}$, de unde $n = 5$);
  \item $x_n = \frac{1}{n^4 (2n)!}, \varepsilon = 10^{-6}$.
\end{enumerate}


\end{document}

